\documentclass[11pt]{article}

\usepackage[margin=1in]{geometry}
\usepackage{amsmath,amssymb,amsthm,mathtools}
\usepackage{booktabs}
\usepackage{graphicx}
\usepackage{microtype}
\usepackage{natbib}
\usepackage[hidelinks]{hyperref}

\newtheorem{assumption}{Assumption}
\newtheorem{lemma}{Lemma}
\newtheorem{proposition}{Proposition}
\newtheorem{corollary}{Corollary}

\title{Working Paper Title}
\author{Author Name}
\date{\today}

\begin{document}

\maketitle

\begin{abstract}
\textbf{Problem definition:} State the operational evolution, focal decision, and unresolved tension.

\textbf{Methodology/results:} State the model class, endogenous choices, two to four linked verified result groups, and the distinctive mechanism, policy comparison, or reversal.

\textbf{Managerial implications:} State the decision maker, design lever, double-edged effect, and relevant boundary.
\end{abstract}

\noindent\textbf{Keywords:} analytical operations management; decision problem; mechanism.

\section{Introduction}
Establish the institutional evolution and first-order benefit, expose the analytical puzzle, state the connected research-question ladder and model, present two to four finding groups aligned with the analysis, distinguish or recover the closest model, and synthesize the bounded implication.

\section{Literature Review}
Organize two to four streams by decision problem and mechanism. End each stream with a consequential difference and close with a cross-stream synthesis.

\section{Model Setup}
Map the institution, then define timing, actors, information, decisions, continuation payoffs, objectives, constraints, equilibrium concept, and notation.

\subsection{Operational Stage}
Define the base-stage decisions and uncertainty.

\subsection{Continuation Environment}
Define downstream or linked payoffs and explain which part depends on the base-stage outcome.

\subsection{Game and Equilibrium Concept}
State beliefs, incentive constraints, deviations, and the solution concept.

\begin{table}[t]
\centering
\caption{Notation}
\begin{tabular}{ll}
\toprule
Symbol & Definition \\
\midrule
$x$ & Focal decision \\
$\theta$ & Primitive or type \\
\bottomrule
\end{tabular}
\end{table}

\begin{lemma}[Complete-information benchmark]\label{lem:benchmark}
Characterize the benchmark contract or decision and expected payoff.
\end{lemma}

Explain why each benchmark choice is optimal and what economic object it maximizes or protects.

\section{Equilibrium Analysis}

\subsection{Mechanism-Removed Benchmark}
Remove the focal continuation value, information linkage, or endogenous feature and characterize the resulting equilibrium.

\subsection{First Enriched Environment}
Add one feature, state the changed payoff or constraint, and characterize the resulting regimes.

\begin{proposition}\label{prop:main}
State one verified result with all required conditions.
\end{proposition}

Translate the equilibrium behavior, compare it with the benchmark, identify the binding constraint and opposing marginal forces, explain any figure or regime switch, and distinguish the closest mechanism in continuous prose.

\subsection{Main Environment with Endogenous Linkage}
Add the focal linkage, characterize any new interior, hybrid, reversal, or selection result, and explain why it cannot arise in the preceding environment.

\subsection{Comparative Statics and Equilibrium Selection}
Separate within-regime effects, discrete switches, and coexistence or selection.

\section{Extensions and Robustness}
Open by naming the headline mechanism and the threats being tested. Use targeted relaxations and classify each result as mechanism-preserving, attenuating, reversing, or replacing.

\section{Conclusion}
Retrace the benchmark-to-main progression, verified mechanism, reversal or selection result, bounded decision implication, and limitations. Introduce no new result.

% Add \bibliographystyle{apalike} and \bibliography{references} after creating references.bib.

\appendix
\section{Proofs and Additional Derivations}

\begin{proof}[Proof of Lemma~\ref{lem:benchmark}]
Derive candidates, verify feasibility and optimality, compare boundaries, and conclude.
\end{proof}

\begin{proof}[Proof of Proposition~\ref{prop:main}]
Partition the parameter space, derive each candidate, compare boundaries and deviations, verify threshold ordering, and conclude.
\end{proof}

\end{document}
